\section*{अध्याय \ref{ch:b3:membrane-traffic} --- झिल्ली यातायात और प्रोटीन छँटाई}

\begin{solution}{exo:b3:membrane-traffic:1}
अमीनो-सिरे का जलविरोधी खंड: \omterm{def:b3:membrane-traffic:signals}{संकेत पेप्टाइड}, अंतर्द्रव्यी जालिका में
(और चूक से आगे पृष्ठ तक)। PKKKRKV:
\omterm{def:b3:membrane-traffic:signals}{केंद्रक-स्थानिकीकरण संकेत}, छिद्र से होकर केंद्रक में।
अमीनो-सिरे की उभयरागी कुंडली: सूत्रकणिकीय \omterm{def:b3:membrane-traffic:signals}{पूर्वानुक्रम}, सूत्रकणिका के
आधात्री में। KDEL: गॉल्जी से जालिका में वापसी।
मैनोज़-6-फ़ॉस्फ़ेट: ट्रांस-गॉल्जी से \omterm{def:b3:membrane-traffic:endocytosis}{अंतःकाय}
और \omterm{def:b3:membrane-traffic:endocytosis}{लयनकाय} तक।
\end{solution}

\begin{solution}{exo:b3:membrane-traffic:2}
(1) बिना झिल्लियों के अनूदित शृंखला स्रावित प्रोटीन से लगभग बीस अवशेष
लंबी थी: स्राव के दौरान कोई खंड हटाया जाता है। (2) अनुवादन के दौरान
झिल्लियाँ उपस्थित हों तो उत्पाद परिपक्व आकार का था और प्रोटिएज़ से
सुरक्षित था: वह पुटिकाओं में घुस चुका था और उनके भीतर खंड खो चुका था।
(3) अनुवादन के बाद जोड़ी गई झिल्लियों ने इनमें से कुछ नहीं किया: प्रवेश
संश्लेषण से युग्मित है --- सह-अनुवादनी --- और वह खंड, यानी \omterm{def:b3:membrane-traffic:signals}{संकेत
पेप्टाइड}, केवल किसी नवजात शृंखला पर ही काम करता है।
\end{solution}

\begin{solution}{exo:b3:membrane-traffic:3}
जालिका से गॉल्जी: \omterm{def:b3:membrane-traffic:coats}{COPII}, अग्रगामी। गॉल्जी से जालिका: \omterm{def:b3:membrane-traffic:coats}{COPI},
प्रतिगामी। ट्रांस-गॉल्जी से \omterm{def:b3:membrane-traffic:endocytosis}{अंतःकाय}: अपने संयोजकों सहित \omterm{def:b3:membrane-traffic:coats}{क्लैथ्रिन},
\omterm{def:b3:membrane-traffic:endocytosis}{लयनकाय} की ओर अग्रगामी। जीवद्रव्य-कला से \omterm{def:b3:membrane-traffic:endocytosis}{अंतःकाय}: \omterm{def:b3:membrane-traffic:coats}{क्लैथ्रिन},
भीतर की ओर (अंतर्ग्रहण)।
\end{solution}

\begin{solution}{exo:b3:membrane-traffic:4}
शर्कराएँ केवल जालिका और गॉल्जी की अवकाशिका के भीतर जुड़ती हैं। हर
पुटिका का अवकाशिकी फलक मुकुलन और संलयन के पार अवकाशिकी ही रहता है, और
जब कोई पुटिका जीवद्रव्य-कला से संलयित होती है तो उसका अवकाशिकी फलक
कोशिका का बाहर बन जाता है; दिशा कभी उलटी नहीं होती, इसलिए जो भीतर
ग्लाइकोसिलित हुआ वह बाहर आ जाता है।
\end{solution}

\begin{solution}{exo:b3:membrane-traffic:5}
$t^{*} = \ln(0.05/0.2)/(0.05 - 0.2) = \ln 4/0.15 = \qty{9.2}{min}$;
$B(t^{*}) = (0.2/(-0.15))(e^{-1.85} - e^{-0.46}) = 0.63$। \qty{60}{min} पर
$A \approx 0$, $B = 1.33\,(e^{-3} - e^{-12}) = 0.066$, इसलिए $C = 0.93$।
\end{solution}

\begin{solution}{exo:b3:membrane-traffic:6}
\omterm{def:b3:membrane-traffic:signals}{संकेत पेप्टाइड} अमीनो सिरे को अवकाशिका में डाल देता है; हर जलविरोधी खंड
बारी-बारी से विराम-स्थानांतरण या आरंभ-स्थानांतरण है, इसलिए चार
झिल्ली-पारी कुंडलियाँ बनती हैं (लगभग अवशेष 30--52,
90--112, 150--172, 210--232)। खंड: अमीनो सिरा अवकाशिकी;
52--90 कोशिकाद्रव्यी; 112--150 अवकाशिकी; 172--210 कोशिकाद्रव्यी;
कार्बोक्सी सिरा अवकाशिकी। ग्लाइकोसिलन केवल अवकाशिकी खंडों पर ---
अमीनो सिरा, 112--150 और कार्बोक्सी सिरा --- जो पृष्ठ पर बाह्यकोशिकीय
बन जाते हैं।
\end{solution}

\begin{solution}{exo:b3:membrane-traffic:7}
$J_{\max} = \num{30000}\times 0.25\times 0.05/0.30 = 1250$ प्रति मिनट;
पृष्ठ-अंश $k_{r}/(k_{i} + k_{r}) = 0.05/0.30 = 0.17$। $R$ दोगुना करने
से $J$ ठीक दोगुना हो जाता है; वापसी-दर ही अड़चन है (अधिकांश ग्राही
भीतर हैं), और $k_{r}$ दोगुना करने से $\num{30000}\times 0.25
\times 0.1/0.35 \approx 2140$ मिलता, यानी $1.7$ का गुणक।
\end{solution}

\begin{solution}{exo:b3:membrane-traffic:8}
(क) ग्राही नहीं: अंकित जल-अपघटनी चूक-मार्ग लेकर स्रावित हो जाते हैं;
\omterm{def:b3:membrane-traffic:endocytosis}{लयनकायों} में एंज़ाइम नहीं रहते और वे अपचित सामग्री से भर जाते हैं ---
कोई संचय-रोग। (ख) फ़ॉस्फ़ोट्रांसफ़रेज़ नहीं: उसी लक्षण-प्ररूप का
दूसरा कारण (I-कोशिका रोग), एंज़ाइम अनंकित और स्रावित। (ग) उदासीन
\omterm{def:b3:membrane-traffic:endocytosis}{अंतःकाय}: ग्राही अपना भार छोड़ ही नहीं सकता, इसलिए वह पुनर्चक्रित नहीं
होता और जल-अपघटनी गलत पहुँचते या स्रावित हो जाते हैं; LDL तथा दूसरे
ग्राही भी चक्कर लगाना बंद कर देते हैं।
\end{solution}

\begin{solution}{exo:b3:membrane-traffic:9}
वह टायरोसीन पुच्छ के उसी \omterm{def:b3:bioinformatics:motif}{अभिप्ररूप} का अंग है जिसे \omterm{def:b3:membrane-traffic:coats}{क्लैथ्रिन} संयोजक
पहचानता है; उसके बिना ग्राही आवरित गड्ढों में बटोरा नहीं जाता।
ग्राही LDL सामान्य रूप से बाँधते हैं पर पृष्ठ पर एक जैसे फैले रहते
हैं, गड्ढों से बाहर, और लिगंड बाहर ही रह जाता है --- यह छँटाई का
रोग है, बंधन का नहीं।
\end{solution}

\begin{solution}{exo:b3:membrane-traffic:10}
$k_{1} = k_{2} = k$ के साथ $B = kt\,e^{-kt}$ आज़माइए: $\mathrm{d}B/\mathrm{d}t
= k e^{-kt} - k^{2}t e^{-kt} = kA - kB$, जैसा चाहिए, और $B(0) = 0$ भी।
शिखर वहाँ जहाँ $1 - kt = 0$: $t^{*} = 1/k$, $B = e^{-1} = 0.37$। सामान्य
रूप में $C(t) = 1 - (k_{2}e^{-k_{1}t} - k_{1}e^{-k_{2}t})/(k_{2} - k_{1})$,
जो $k_{1}$ और $k_{2}$ की अदला-बदली में सममित है: अकेला कणिका-वक्र
नहीं बता सकता कि दोनों कक्षों में तेज़ कौन-सा है।
\end{solution}

\begin{solution}{exo:b3:membrane-traffic:11}
हर पुटिका का क्षेत्रफल $4\pi(50\,\text{nm})^{2} = \qty{0.031}{\micro
m^2}$ है; प्रति मिनट एक हज़ार यानी \qty{31}{\micro m^2/min}, इसलिए पूरा
पृष्ठ $3000/31 \approx \qty{95}{min}$ में। कोशिका इसलिए नहीं सिकुड़ती
कि झिल्ली बहिर्ग्रहण (पुनर्चक्रण पुटिकाओं) से उसी दर पर लौट आती है
--- प्रवाहों का संतुलन। हर पुटिका $5\times 10^{-19}$ L रखती है; प्रति
मिनट एक हज़ार, एक घंटे के लिए, $3\times 10^{-14}$ L हुआ, यानी किसी
\qty{2}{pL} कोशिका के आयतन का प्रति घंटा लगभग \qty{1.5}{\%}।
\end{solution}

\begin{solution}{exo:b3:membrane-traffic:12}
बोटुलिनम विष चालक सिरे में काम करता है: कोई ऐसीटिलकोलीन मुक्त नहीं
होती, पेशी को कोई आदेश नहीं मिलता और वह ढीली पड़ी रहती है --- शिथिल
पक्षाघात। धनुस्तंभ विष चालक तंत्रिकाक्ष के सहारे ऊपर और उन निरोधी
अंतर्न्यूरॉनों तक पहुँचाया जाता है जो सामान्यतः चालक न्यूरॉनों को थामे
रखते हैं; उनका \omterm{def:b3:membrane-traffic:fusion}{SNARE} कट जाने पर कोई ग्लाइसीन या GABA मुक्त नहीं होती,
चालक न्यूरॉन बेरोक चलते हैं और हर पेशी सिकुड़ जाती है --- स्पास्टिक
पक्षाघात। किसी अति-सक्रिय पेशी में अंतःक्षेपित बोटुलिनम विष के कुछ
नैनोग्राम उसे महीनों तक स्थानिक रूप से चुप कर देते हैं (जब तक नया
\omterm{def:b3:membrane-traffic:fusion}{SNARE} न बने और सिरा उबर न जाए): दुस्तानताओं, स्पास्टिकता, भेंगापन और
झुर्रियों का एक उपचार।
\end{solution}

\begin{solution}{pb:b3:membrane-traffic:1}
\textbf{1.} $10^{7}\times \qty{25000}{Da}\times 1.66\times 10^{-24}$ g
$= \qty{0.42}{pg}$ प्रति मिनट, \qty{600}{pg} प्रति दिन --- अपनी
प्रोटीन-मात्रा से दोगुना।
\textbf{2.} पुटिका का आयतन $\tfrac{4}{3}\pi\,35^{3} = 1.8\times 10^{5}$
nm$^{3}$, उसका \qty{30}{\%} $5.4\times 10^{4}$; एक एंज़ाइम अणु
$\tfrac{4}{3}\pi\,2^{3} = 34$ nm$^{3}$: प्रति पुटिका लगभग \num{1600}
अणु; प्रति मिनट $10^{7}/1600 \approx 6200$ पुटिकाएँ।
\textbf{3.} प्रति पुटिका क्षेत्रफल $4\pi\,35^{2} = \qty{0.0154}{\micro m^2}$, प्रति मिनट \qty{95}{\micro
m^2}: जालिका का \qty{6000}{\micro m^2} लगभग एक घंटे में, यदि कुछ भी
वापस न आए।
\textbf{4.} $\qty{95}{\micro m^2} = 9.5\times 10^{7}$ nm$^{2}$, दो परतें
\qty{0.7}{nm^2} पर: प्रति मिनट $2.7\times 10^{8}$ लिपिड। कोशिका \omterm{def:b3:membrane-traffic:coats}{COPI}
पुटिकाओं और पुनर्चक्रण से झिल्ली लौटाती है, और जो कणिकाओं में सदा के
लिए चली जाती है उसकी भरपाई के लिए लिपिड संश्लेषित करती है।
\textbf{5.} कणिका का आयतन $5.2\times 10^{8}$ nm$^{3}$, उसका \qty{30}{\%}
$1.6\times 10^{8}$: प्रति कणिका $4.7\times 10^{6}$ अणु; प्रति घंटा $6\times
10^{8}$ अणु लगभग $130$ कणिकाएँ भर देते हैं।
\textbf{6.} हर कणिका $\pi d^{2} = \qty{3.1}{\micro m^2}$ जोड़ती है; $300$
कणिकाएँ \qty{50}{\micro m^2} में \qty{940}{\micro m^2} जोड़ती हैं:
बीस गुना। झिल्ली को मिनटों में प्रतिपूरक अंतर्ग्रहण से वापस लेना पड़ता है।
\textbf{7.} $t^{*} = \ln 2/0.05 = \qty{13.9}{min}$, $B = 0.50$।
\textbf{8.} $A = e^{-3} = 0.05$; $B = 2(e^{-1.5} - e^{-3}) = 0.35$;
$C = 0.60$।
\textbf{9.} \qty{10}{min} और \qty{20}{min}; माध्य कुल \qty{30}{min}।
\textbf{10.} $t^{*} = \ln(0.1)/(-0.09) = \qty{25.6}{min}$; $B(t^{*}) =
(0.1/(-0.09))(e^{-2.56} - e^{-0.256}) = 0.77$: गॉल्जी फूल जाता है, भार
के तीन-चौथाई से भरा हुआ।
\textbf{11.} शिखर पर $k_{1}e^{-k_{1}t^{*}} = k_{2}e^{-k_{2}t^{*}}$;
इसे $B$ में रखने पर $B_{\max} = (k_{1}/k_{2})^{\,k_{2}/(k_{2} -
k_{1})}$ मिलता है: घात $k_{2}/(k_{2} - k_{1})$ है। $k_{1} > k_{2}$ के
लिए घातांक ऋणात्मक है और आधार $1$ से ऊपर; $k_{1} < k_{2}$ के लिए
आधार $1$ से नीचे है और घातांक धनात्मक --- दोनों स्थितियों में मान
$1$ से नीचे रहता है (जाँच: यहाँ $2^{-1} = 0.5$)।
\textbf{12.} $k_{1} \gg k_{2}$ के लिए $B(t) \to e^{-k_{2}t}$: चिह्न
तुरंत गॉल्जी में होता है और $k_{2}$ से निकलता है --- जालिका का चरण
अदृश्य है और गॉल्जी का वक्र एक सरल क्षय है।
\textbf{13.} $\theta = 2/(2+2) = 0.5$; $J = \num{50000}\times 0.2\times
0.1\times 0.5/(0.1 + 0.1) = 2500$ कण प्रति मिनट, प्रति घंटा \num{150000}।
\textbf{14.} प्रति घंटा $1.5\times 10^{5}\times 1500 = 2.3\times 10^{8}$
कोलेस्टेरॉल अणु; $5\times 10^{9}/2.3\times 10^{8} \approx 22$ घंटे।
\textbf{15.} पृष्ठ-अंश $k_{r}/(k_{r} + k_{i}\theta) = 0.5$; एक
चक्कर $1/(k_{i}\theta) + 1/k_{r} = 10 + 10 = \qty{20}{min}$ लेता है:
बीस घंटों में लगभग $60$ चक्कर।
\textbf{16.} $J = 1250$ प्रति मिनट। संतृप्ति से नीचे $J \propto R\,[L]$,
इसलिए आधे ग्राहियों के साथ उसी निकासी के लिए प्लाज़्मा LDL दोगुना होना चाहिए।
\textbf{17.} $R$ फिर \num{50000}: निकासी लौट आती है, प्लाज़्मा LDL
सामान्य की ओर गिर जाता है --- स्टैटिन का प्रभाव।
\textbf{18.} बिना ग्राहियों के LDL केवल धीमे अविशिष्ट मार्गों से साफ़
होता है, इसलिए वह घंटों के बजाय दिनों तक घूमता है, ऑक्सीकृत हो जाता है,
और महाभक्षकों के मार्जक ग्राही उसे ग्रहण कर लेते हैं, जो धमनीकाठिन्य
की पट्टिकाओं की फेन-कोशिकाएँ बन जाते हैं; और स्टैटिन, जो ग्राही
प्रेरित करके काम करते हैं, के पास प्रेरित करने को कुछ नहीं होता।
\textbf{19.} $V = \tfrac{4}{3}\pi(0.25)^{3} = \qty{0.065}{\micro m^3} =
6.5\times 10^{-17}$ L। pH $7.2$ पर: $6.3\times 10^{-8}\times 6.5\times
10^{-17} = 4\times 10^{-24}$ mol --- लगभग $2$ मुक्त प्रोटॉन। pH $4.7$ पर:
$2\times 10^{-5}\times 6.5\times 10^{-17} = 1.3\times 10^{-21}$ mol,
लगभग $780$।
\textbf{20.} $50\times 780 \approx \num{39000}$ प्रोटॉन; $50$ पंप,
प्रति सेकंड $200$ की दर से, प्रति सेकंड \num{10000} चलाते हैं: लगभग \qty{4}{s}।
\textbf{21.} धन आवेश भीतर पंप करने से अवकाशिका धनात्मक हो जाती है, और
कुछ हज़ार आवेश पंप रोक देते; आवेश उदासीन करने के लिए क्लोराइड किसी
वाहिका से भीतर आता है (या धनायन बाहर जाते हैं)।
\textbf{22.} उदासीन pH पर निष्क्रिय एंज़ाइम कोई हानि नहीं करते यदि कोई
\omterm{def:b3:membrane-traffic:endocytosis}{लयनकाय} रिसे या यदि वे गलती से कोशिकाद्रव्य या पृष्ठ पर पहुँच जाएँ:
अम्ल की आवश्यकता पाचन को उसी कक्ष तक सीमित रखती है जो उसके लिए बना है।
\textbf{23.} दोनों में कोई pH संवेदक होना चाहिए --- हिस्टिडीन, जिनका
pH~6 के पास प्रोटॉनन बंधन-पृष्ठ बदल देता है: LDL ग्राही प्रोटॉनित होने
पर अपने ही बंधन-स्थल पर कोई प्रांत मोड़ देता है, और
मैनोज़-6-फ़ॉस्फ़ेट ग्राही भी इसी तरह आसक्ति खो देता है।
\textbf{24.} $10^{7.2 - 4.7} = 10^{2.5} \approx 300$ गुना सांद्रण;
फँसा हुआ क्षार प्रोटॉन खपाता है, लयनकायी pH उठा देता है, और
जल-अपघटनियों को निष्क्रिय कर देता है --- इसी तरह क्लोरोक्वीन मलेरिया
के परजीवियों को उनकी भोजन-रसधानी में मारती है और इसी तरह वह प्रयोगों
में \omterm{def:b3:membrane-traffic:endocytosis}{स्वभक्षण} रोकती है।
\textbf{25.} गॉल्जी का शिखर \qty{14}{min} पर; प्रति मिनट कोई \num{6200}
पुटिकाएँ जालिका से निकलती हैं; प्रति घंटा \num{150000} LDL कण ग्रहण
किए जाते हैं।
\end{solution}
